\bookchapter{נספח לפרק 23: מדריך למתחילים ל-GitHub – מאחסון קוד לפלטפורמת שיתוף פעולה}

\chapterintro{בפרק הקודם דיברנו על “חשיבה תכנותית”, וכדי באמת ליישם חשיבה תכנותית ולעבוד על פרויקטים עם אחרים, אי אפשר להתעלם משם אחד – GitHub. בין אם אתה סטודנט למדעי המחשב שרוצה להשתתף בתחרויות, רוצה לשמור את יצירות הקוד שלך, או פשוט רוצה לראות במה מתעסקים מפתחים מכל העולם, GitHub היא פלטפורמה שאתה חייב להכיר. בחלק זה נתחיל מהמושגים הבסיסיים ביותר, ונבנה הבנה ברורה.}

\booksection{מה זה GitHub בעצם?}

המבוססת על Git. היא עלתה לאוויר ב-2008, נרכשה על ידי מיקרוסופט ב-2018, ועד היום היא פלטפורמת אחסון הקוד הפתוח הגדולה בעולם, מאחסנת מאות מיליוני מאגרים, ומכסה כמעט את כל שפות התכנות המרכזיות.\textbf{אם נשתמש באנלוגיה להבנה:}Git הוא יומן במחשב שלך, GitHub היא פלטפורמת בלוגים ציבורית

– אתה יכול לפרסם את היצירות שלך, לעקוב אחר פרויקטים של אחרים, להשתתף בדיונים ולשתף פעולה. היא לא רק “מחסן אחורי למתכנתים”, אלא יותר כמו “עיר טכנולוגית” שמפתחים מכל העולם בונים יחד – פרויקטים נולדים כאן, עוברים איטרציות, נדונים, ומשומשים מחדש.\textbf{עבור אנשים רגילים, GitHub גם אינה רחוקה מהישג יד. הרבה לא-מתכנתים משתמשים בה לאחסון הערות, ספרים אלקטרוניים, בלוגים אישיים, ואפילו לגיבוי קבצי תצורה. כל עוד יש לך ולו מעט צורך “לארגן את הדברים שלך ולשתף אותם”, GitHub שווה שתכיר.}2. ההבדל בין Git ל-GitHub (מושג מפתח!)

הרבה אנשים מתבלבלים בין שני המונחים האלה בפעם הראשונה, הנה הסבר ברור:

\booksection{ההבדל בין Git ל-GitHub (מושג מפתח!)}

Git

\begin{center}
\begin{tabularx}{\linewidth}{|X|X|X|}
\hline
\textbf{GitHub} & \textbf{מהות} & \textbf{כלי בקרת גרסאות (תוכנה)} \tabularnewline
\hline
פלטפורמת אחסון קוד (אתר) & מיקום ריצה & על המחשב המקומי \tabularnewline
בענן (גישה דרך דפדפן) & אחריות ליבה & תיעוד כל שינוי בקובץ (איך שינו, מי שינה, מתי שינו) \tabularnewline
העלאת מאגרי Git לענן, לנוחות אחסון, הצגה ושיתוף פעולה & האם צריך רשת & לא צריך, זמין במצב לא מקוון \tabularnewline
צריך חיבור לאינטרנט & שייכות & תוכנת קוד פתוח, נוצרה על ידי Linus Torvalds (אבי Linux) \tabularnewline
מוצר בבעלות מיקרוסופט & מתחרים דומים & SVN, Mercurial \tabularnewline
GitLab, Gitee (Gitee), Bitbucket & בקיצור: & Git מנהל "איך לתעד שינויים", GitHub מנהל "איפה השינויים נמצאים ואיך לשתף פעולה" \tabularnewline
\hline
\end{tabularx}
\end{center}

. Git הוא כלי, שרץ מקומית; GitHub היא פלטפורמה, שמעלה מאגרי Git לענן, לנוחות שיתוף פעולה והצגה. אתה יכול להשתמש רק ב-Git בלי GitHub (למשל לכתוב קוד לבד עם ניהול גרסאות), אבל כדי להשתמש ב-GitHub חייבים להשתמש ב-Git (כי GitHub מבוססת על Git).\textbf{3. מושגי ליבה ופונקציות של GitHub}מושגי יסוד

\booksection{מושגי ליבה ופונקציות של GitHub}

\booksubsection{מושגי יסוד}

: פרויקט אחד מתאים למאגר אחד, שבתוכו מאוחסנים כל קבצי הקוד, הגרסאות ההיסטוריות והתיעוד של הפרויקט. אפשר להבין אותו כ“תיקיית פרויקט”.

\begin{itemize}
  \item \textbf{הגשה (Commit)}: תיעוד של שינוי בקוד. כל Commit מגיע עם הודעת הסבר, המתעדת "מה שיניתי הפעם ולמה". זהו הבסיס ליכולת "מסע בזמן" של Git/GitHub.
  \item \textbf{ענף (Branch)}: קו פיתוח עצמאי שמתפצל מהקו הראשי. לדוגמה, אם אתה רוצה להוסיף פונקציה חדשה, אתה יכול לפתוח ענף ולפתח בנפרד, ולאחר סיום למזג חזרה לקו הראשי, בלי להשפיע על אחרים.
  \item \textbf{שכפול (Clone)}: העתקה מלאה של מאגר מהענן למחשב המקומי, שלאחריה ניתן לערוך קוד מקומית.
  \item \textbf{דחיפה (Push) / משיכה (Pull)}: Push הוא העלאת קוד מקומי ל-GitHub, Pull הוא הורדת הקוד העדכני מ-GitHub למחשב המקומי. אלה שתי פעולות הליבה לסנכרון קוד בעבודת צוות.
  \item \textbf{בקשת מיזוג (Pull Request, בקיצור PR)}: כשסיימת לערוך קוד בענף ורוצה למזג לענף הראשי, אתה מגיש PR, ומאפשר לחברי הצוות לסקור את הקוד לפני המיזוג. זהו מנגנון הליבה של שיתוף הפעולה ב-GitHub.
  \item \textbf{נושא (Issue)}: משמש לתיעוד באגים, דרישות פונקציונליות, נושאי דיון. כל Issue הוא כמו פוסט דיון, ניתן להקצות אחראי ולהוסיף תגיות.
  \item \textbf{Fork}: העתקת מאגר של מישהו אחר לחשבון שלך. אתה יכול לשנות כרצונך בעותק שלך, מבלי להשפיע על המאגר המקורי. לאחר השינויים תוכל לשלוח PR למחבר המקורי, לבקש למזג את השיפורים שלך.
  \item \textbf{פונקציות ליבה}GitHub מציעה לא רק אחסון קוד, אלא גם:
\end{itemize}

\booksubsection{פונקציות ליבה}

: תיעוד ההיסטוריה של כל שינוי בקוד, אם קלקלת קוד תמיד אפשר לחזור לגרסה התקינה הקודמת.

\begin{itemize}
  \item \textbf{שיתוף פעולה צוותי}: תמיכה בפיתוח מקבילי של מספר משתמשים באותו פרויקט, פיתוח בו-זמני של מספר אנשים ללא דריסה וללא התנגשויות.
  \item \textbf{מעקב בעיות}: שימוש ב-Issue לתיעוד באגים ובקשות פונקציונליות, נוח למעקב צוותי.
  \item \textbf{ניהול פרויקטים}: כלים משולבים לניהול משימות, לוחות פרויקטים (Projects) וכו', שווה ערך ל-Trello קל משקל מובנה.
  \item \textbf{אוטומציה}: באמצעות
  \item \textbf{אוטומציה}ניתן להשיג קומפילציה אוטומטית, בדיקות, אריזה ופריסה לאחר הגשת קוד, שזהו כלי הליבה של "CI/CD (אינטגרציה רציפה / פריסה רציפה)" המודרני.\textbf{ניהול תיעוד}: README מציג את הפרויקט, Wiki מאחסן ארכיטקטורה מפורטת ותקני פיתוח, Pages אפילו יכול לארח אתרים סטטיים ישירות.
  \item \textbf{אחסון קוד והצגה}: לכל מאגר יש תצוגת דף אינטרנט יפה, תומכת בהדגשת תחביר, ה"חלון ראווה" של יצירותיך הטכניות.
  \item \textbf{קהילת קוד פתוח}: אתה יכול לתת Star (מועדפים) לפרויקטים שאתה אוהב, Watch (מעקב) אחר עדכוני פרויקטים, Follow (מעקב) אחר מפתחים, ולהשתתף בשיתוף פעולה גלובלי בקוד פתוח.
  \item \textbf{4. הרשמה והתחלה: צור את המאגר הראשון שלך בחמש דקות}הרשמת חשבון
\end{itemize}

\booksection{הרשמה והתחלה: צור את המאגר הראשון שלך בחמש דקות}

\booksubsection{הרשמת חשבון}

\begin{enumerate}
  \item , לחץ "Sign up".\href{https://github.com}{מלא אימייל, הגדר סיסמה ושם משתמש (שם המשתמש יופיע בכתובת הדף האישי שלך, מומלץ באנגלית, לדוגמה}github.com/yourname
  \item ).\texttt{השלם אימות אימייל ואימות אנושי.\allowbreak }בחר תוכנית חינמית (Free), לגמרי מספיקה למשתמש אישי.
  \item יצירת מאגר ראשון
  \item לאחר התחברות, לחץ על סימן "+" בפינה הימנית העליונה, בחר "New repository".
\end{enumerate}

\booksubsection{יצירת מאגר ראשון}

\begin{enumerate}
  \item my-first-project
  \item כתבו תיאור קצר (Description), כדי שאחרים יידעו במה עוסק המאגר.
  \item בחרו ציבורי (Public) או פרטי (Private). מאגר ציבורי גלוי לכולם, פרטי גלוי רק לכם ולמשתפי הפעולה. גם חשבון חינמי תומך במאגרים פרטיים ללא הגבלה.
  \item סמן "Add a README file" – README הוא "מדריך ההוראות" של הפרויקט, מומלץ לכל מאגר.
  \item לחץ "Create repository", ומאגר GitHub הראשון שלך נולד!
  \item טיפ קטן: מה זה README
  \item README.md הוא קובץ בפורמט Markdown, זהו התוכן הראשון שאחרים רואים כשהם פותחים את המאגר שלך. README טוב כולל בדרך כלל: סקירת פרויקט, שיטת התקנה, דוגמאות שימוש, רישיון, פרטי יצירת קשר. כתיבת README טוב היא מיומנות בסיסית של מתכנת, והיא הסף הראשון לכך שפרויקט יהיה ניתן לשימוש על ידי אחרים.
\end{enumerate}

\begin{tcbinfo}{5. זרימת עבודה בסיסית: ממקומי לענן}
כדי להבין את זרימת העבודה של GitHub, המפתח הוא לזכור את המחזור הזה:

\end{tcbinfo}

\booksection{זרימת עבודה בסיסית: ממקומי לענן}

. הנה הסבר במילים פשוטות של כל שלב:\textbf{עריכת קבצים}: כתוב קוד או ערוך מסמכים בעורך במחשב המקומי.

\begin{enumerate}
  \item \textbf{git add}: הוסף את הקבצים ששונו ל"אזור הזמני", שווה ערך לאמירה ל-Git "אני מתכונן לתעד את השינויים האלה".
  \item \textbf{git commit}: תעד רשמית את השינויים באזור הזמני כ"הגשה", צרף הסבר, לדוגמה "תיקון באג התחברות". שלב זה מתבצע רק מקומית.
  \item \textbf{git push}: העלה את ההגשות המקומיות לענן GitHub, כך שאחרים יוכלו לראות.
  \item \textbf{git pull}: לפני תחילת עבודה, משוך את הקוד העדכני מהענן, וודא סנכרון מקומי עם שאר הצוות.
  \item \textbf{למתחילים, זכור את הנוסחה הזו:}"קודם pull, אחר כך ערוך, סיים לערוך add, commit עם הסבר, ולבסוף push למעלה"
\end{enumerate}

. תאמץ את ההרגל הזה, ושיתוף פעולה צוותי לא יתבלגן.\textbf{6. זרימת עבודה לשיתוף פעולה צוותי: איך Pull Request עובד}החלק החזק ביותר של GitHub הוא מנגנון שיתוף הפעולה שלה. זרימת עבודה טיפוסית לשיתוף פעולה צוותי נראית כך:

\booksection{זרימת עבודה לשיתוף פעולה צוותי: איך Pull Request עובד}

: מהענף הראשי (בדרך כלל נקרא main או master) פתח ענף חדש, לדוגמה

\begin{enumerate}
  \item \textbf{feature/login}, ופתח את הפונקציה החדשה בענף זה.\texttt{פתח והגש}: כתוב קוד בענף החדש, בצע מספר commit לתיעוד התקדמות.
  \item \textbf{דחוף ל-GitHub}: דחוף את הענף למאגר בענן.
  \item \textbf{הגש Pull Request}: בדף GitHub הגש PR, בקש למזג את הענף שלך לענף הראשי. כתוב בבירור אילו שינויים בוצעו ב-PR ולמה.
  \item \textbf{סקירת קוד (Code Review)}: חברי הצוות יבדקו את הקוד שלך, יעירו הערות או יציעו שינויים. אתה יכול להמשיך לבצע commit לפי המשוב, ה-PR יתעדכן אוטומטית.
  \item \textbf{מיזוג}: לאחר אישור הסקירה, לחץ "Merge" למזג את הקוד לענף הראשי. התרומה שלך נכנסה רשמית לפרויקט!
  \item \textbf{זרימת עבודה זו מבטיחה}שכל קוד שנכנס לענף הראשי עובר סקירה
\end{enumerate}

, מה שמפחית משמעותית את הסיכוי לשגיאות. פרויקטי קוד פתוח משתמשים באותו מנגנון לקבלת תרומות מזרים מכל העולם – אתה לגמרי יכול להגיש PR לפרויקט קוד פתוח מוכר, ואם הוא מתקבל אתה הופך לתורם של הפרויקט.\textbf{7. תרחישי שימוש טיפוסיים של GitHub}סטודנטים המשתתפים בתחרויות ופרויקטי קורס

\booksection{תרחישי שימוש טיפוסיים של GitHub}

\booksubsection{סטודנטים המשתתפים בתחרויות ופרויקטי קורס}

גם לא-מתכנתים יכולים להשתמש ב-GitHub. אנשים רבים משתמשים בה לאחסון הערות לימוד, בלוגים טכניים (בשילוב GitHub Pages ניתן לבנות אתר בחינם), ספרים אלקטרוניים. כתיבה בפורמט Markdown + ניהול גרסאות Git + גיבוי בענן, שלושה באחד.

\booksubsection{בסיס ידע אישי ואחסון הערות}

GitHub היא הזירה המרכזית של תנועת הקוד הפתוח העולמית. אתה יכול לעשות Fork לפרויקטים שאתה אוהב, לשנות ולשלוח PR חזרה לקהילה; או פשוט לתת Star לשמירת פרויקטים טובים ולעקוב אחר עדכונים. השתתפות בקוד פתוח היא דרך מצוינת לשפר מיומנויות ולצבור מוניטין.

\booksubsection{השתתפות בפרויקטי קוד פתוח}

באמצעות GitHub Actions, ניתן להשיג תהליך אוטומטי מלא של “קוד נדחף, בדיקות אוטומטיות + פריסה אוטומטית לשרת”. הרבה בלוגים אישיים ואתרי תיעוד בנויים כך, כותבים פוסט, עושים push, והאתר מתעדכן אוטומטית.

\booksubsection{פריסה אוטומטית ו-CI/CD}

השתמש במאגר הקוד ככונן ענן, כשאתה מחליף מחשב או עובד ממספר מכשירים, git pull אחד ואתה מקבל את הקוד העדכני, לא צריך יותר להעביר עם דיסק און קי.

\booksubsection{גיבוי קוד וסנכרון רב-מכשירי}

הסבר ההגשה צריך להיות ברור

\booksection{כמה עצות למתחילים}

\begin{itemize}
  \item \textbf{הגש לעתים קרובות, צעדים קטנים ורצים}: כל פעם שמסיימים פונקציה קטנה תעשה commit, במקום לעשות commit אחד רק בסוף יום שלם. ככה קל יותר לאתר בעיות.
  \item \textbf{פתח בענף, אל תשנה ישירות את הענף הראשי}: אפילו בפרויקט אישי, תאמץ הרגל לפתוח ענף, בעתיד כשיהיה שיתוף פעולה צוותי לא תהיה מבולבל.
  \item \textbf{השתמש היטב ב-.gitignore}: קובץ זה אומר ל-Git אילו קבצים לא לכלול בניהול הגרסאות (למשל קבצי תצורה עם סיסמאות, תוצרי קומפילציה, תיקיות תלויות), כדי להימנע מהעלאת קבצי זבל.
  \item \textbf{הסתכל על פרויקטי קוד פתוח מצוינים}: שוטט ב-GitHub Trending, ראה את מבנה הקוד, אופן כתיבת README, ותהליכי שיתוף הפעולה של פרויקטים עם כוכבים רבים, זוהי הדרך החינמית המהירה ביותר להשתפר.
  \item \textbf{בקש חבילת סטודנט של GitHub}: לאחר אימות סטודנט ניתן לקבל בחינם כמות גדולה של כלים ושירותים בתשלום (Copilot, שמות מתחם, שירותי ענן וכו'), בשווי אלפי דולרים.
  \item \textbf{אזהרת אבטחה}לעולם אל תגיש ל-GitHub מידע רגיש כמו סיסמאות, מפתחות API, מחרוזות חיבור למסד נתונים, במיוחד במאגרים ציבוריים – ברגע שדולף מידע, הוא ייסרק על ידי זחלנים תוך שניות, עם השלכות חמורות.
\end{itemize}

\begin{tcbwarn}{אם הגשת בטעות מידע רגיש, לא מספיק למחוק את הקובץ (הוא עדיין בהיסטוריה), צריך להשתמש ב-}
\begin{itemize}
  \item git filter-branch
  \item או בכלים מקצועיים לניקוי ההיסטוריה, ולהחליף מיד את המפתחות שדלפו.\texttt{מאגרים פרטיים גם אינם בטוחים לחלוטין, מידע רגיש עדיף לנהל עם משתני סביבה, שלא ייכנס למאגר הקוד.\allowbreak }הערות לגבי גישה ל-GitHub מסין
  \item הגישה ל-GitHub מסין לעתים אינה יציבה, עלולות להיות בעיות של אי-יכולת לפתוח או שיבוט איטי. ניתן לנסות: (1) להשתמש באתרי מראה סיניים (כמו GitClone, ghproxy) להאצת שיבוט; (2) להגדיר Hosts או להשתמש בפרוקסי; (3) אם משתפים פעולה בעיקר מול קהל סיני, ניתן לשקול גם את התחליף המקומי
\end{itemize}

\end{tcbwarn}

\begin{tcbinfo}{Gitee (Gitee)}
, עם פונקציונליות דומה ומהירות גישה טובה. צוותים רבים מארחים במקביל ב-GitHub וב-Gitee, משלבים בין-לאומיות ומהירות מקומית.\textbf{טריק מתקדם: בניית בלוג אישי חינמי עם GitHub Pages}GitHub Pages מאפשר לך לארח אתר סטטי בחינם, בשילוב עם מחוללי אתרים סטטיים כמו Jekyll, Hexo, Hugo, אפשר לבנות בקלות בלוג אישי. כותבים פוסט Markdown, דוחפים למאגר, הבלוג מתעדכן אוטומטית, עם HTTPS ושם מתחם מותאם אישית מובנים. אינספור בלוגרים טכניים בנו כך את הבלוגים שלהם, אפס עלות, ללא פרסומות, שליטה מלאה בגרסאות.

\end{tcbinfo}

\textbf{9. סיכום: למה כל אחד צריך להכיר את GitHub}\par

GitHub מזמן הפסיקה להיות רק “כלי למתכנתים”, היא יותר כמו

\booksection{סיכום: למה כל אחד צריך להכיר את GitHub}

: ניהול גרסאות הופך הכול לניתן למעקב, מנגנוני שיתוף פעולה מאפשרים לצוותים לעבוד ביעילות, רוח הקוד הפתוח מאפשרת לידע לזרום בחופשיות. גם אם אינך כותב קוד, הבנת המושגים “הגשה, ענף, מיזוג” יכולה לעזור לך לנהל טוב יותר כל תוכן שדורש “עריכה חוזרת + שיתוף פעולה של מספר אנשים” – מתזות ועד מסמכים, מהערות ועד הצעות.\textbf{עבור סטודנטים למדעי המחשב, GitHub היא מיומנות חובה: היא תיק העבודות שלך, פלטפורמת שיתוף הפעולה שלך, מאגר משאבי הלמידה שלך, וגם “קורות החיים” המשכנעים ביותר בחיפוש עבודה עתידי. הירשם מוקדם, העלה את מטלות הקורס והפרויקטים הקטנים שלך, צבור בהדרגה, ובעוד כמה שנים תהיה לך היסטוריה טכנית שתרשים כל מראיין.}בוא נתחיל: השלם את חוויית GitHub הראשונה שלך

הירשם לחשבון GitHub, מלא את הפרופיל האישי (תמונת פרופיל, תיאור קצר, תחומי טכנולוגיה שמעניינים אותך).

\begin{tcbpractice}{צור מאגר ציבורי בשם}
\begin{enumerate}
  \item hello-github
  \item , סמן "Add a README file".\texttt{היכנס למאגר, לחץ על כפתור העריכה של README.\allowbreak md (אייקון עיפרון), כתוב הצגה עצמית ב-\allowbreak Markdown, הגש את השינוי.\allowbreak }בשורת החיפוש של GitHub חפש מילת מפתח שמעניינת אותך (לדוגמה
  \item python-tutorial
  \item ), תן Star לשני פרויקטים שאתה אוהב.\texttt{(מתקדם) התקן Git במחשב, בצע Clone למאגר למחשב המקומי, ערוך את README ואז דחוף חזרה לענן, השלם זרימת עבודה מלאה של "מקומי ענן".\allowbreak }פרק 24: ניתוח מעמיק של הרשאות Root ונעילת BL
  \item פרק זה הוא הרחבה מעמיקה של תוכן פרק 20, מתאים למשתמשים שרוצים באמת להבין את מנגנוני השכבה הנמוכה של הטלפון, ואפילו לנסות להתקין ROM. שים לב, פעולות אלה כרוכות בסיכון מסוים, אנא נהג בזהירות.
\end{enumerate}

\end{tcbpractice}
